\documentclass{ximera}
\title{Inverse Trig CW}

\newcommand{\pskip}{\vskip 0.1 in}

\begin{document}
\begin{abstract}
Applications of inverse trig derivatives.
\end{abstract}
\maketitle

\begin{question} \label{Qpfefd9311}
The ends of a $25$-foot ladder rest against a vertical wall and a floor inclined at an angle of
\[
     \phi = \arccos(5/6)
\]
below the horizontal.

%The ladder is still $25$ feet long and its top end is $20$ feet above the base of the wall.

\begin{onlineOnly}
    \begin{center}
\desmos{q2g6qqwlzc}{900}{600}
\end{center}
\end{onlineOnly}

\href{https://www.desmos.com/calculator/q2g6qqwlzc}{151: Ladder and Sloping Ground}

\begin{enumerate}

\item Find a function that expresses the acute angle the ladder makes with the wall in terms of the distance from the top of the ladder to the base of the wall. Start with a sentence that defines the output, the input, and the function name. Include units in your definition and choose meaningful variable names. Include a picture to help with your explanation.

\item Now suppose the ladder's top end is $20$ feet above the base of the wall.

Find the  scaling factor that converts (by multiplication) a small change in the height of the ladder's top end above the base of the wall (in feet) tp an accurate approximation of the corresponding change in the acute angle the ladder makes with the wall. Include units.

\item Interpret the meaning of the scaling factor with an example using \emph{specific} small changes..



\end{enumerate}



\end{question}



\begin{question} \label{Qodfer3r}
A rocket orbits a planet of radius $2000$km at an altitude of $400$ km above the surface. We wish to find a scaling factor that converts (by multiplication) a small change in orbit's radius to an accurate approximation of the corresponding change in radius of the portion of the planet's surface visible to an astronaut in the rocket. The radius is measured along the planet's surface.

For help see 

\href{https://ximera.osu.edu/calcone/Calculus1/Horizon2/Horizon2}{Looking Down on the Moon Part 2}

\begin{enumerate}

\item Start by writing a complete sentence that defines a function expressing the radius of the visible portion in terms of the altitude above the surface. Use meaningful variable names. Include units in your definition.

\item Use your function to compute the exact value of the scalar factor desribed above.

\item Explain the meaning of the scaling factor by giving an example with specific small changes.

\end{enumerate}

\end{question}

\begin{question} \label{Qdfef3r00} An elastic cord runs from the top of an $80$-foot tall tower to a point on the ground. % $50$ feet from the base of the tower. 

\begin{enumerate}

\item Find a function that expresses the acute angle (in radians)  the cord makes with the ground in terms of the cord's length (in feet). Start with a complete sentence that defines a function with meaningful variable names and units.

\item Suppose the cord is attached to the ground $40$ feet from the base of the tower. Find the scaling factor that converts a small change in the cord's length (in feet) to an accurate approximation of the change in the angle the cord makes with the ground.

\item Explain the meaning of the scaling factor with an example using specific changes.

\item How far from the base of the tower is the end of the cord when the scaling factor in part (b) is $-1/30$ rad/ft?

\end{enumerate}
\end{question}

\begin{question} \label{Wpfee}
The function
\[ 
 h = f(\phi) = \frac{24}{\pi} \arccos \left(-\sqrt{3}\tan\phi   \right)
\]
expresses the number of hours of daylight per day at a latitude of $60^\circ$N in terms of the declination of the sun. The declination is the angle the sun’s rays make with the plane of the equator. It is positive between the spring and fall equinoxes in the northern hemisphere, negative in the fall and winter months.  

\begin{enumerate}
\item Find an expression for the derivative $dh/d\phi$.

\item Evaluate the derivative
\[
     \frac{dh}{d\phi}\Big|_{\phi = \arctan(1/3)} .
\]

\item Interpret the derivative in part (b) as a scaling factor. Give an example with specific small changes.  
\end{enumerate}
\end{question}


\begin{question} \label{Qdefefr}
The function
\[
 h = f(s) \, , \, 0\leq s \leq 10,
\]
expresses the altitude (in km) of a trail in terms of the distance from the trailhead (in km).

\begin{enumerate}
\item Find a function
\[
  \phi = g(\theta) \, , \, -\pi/2 <  \theta < \pi/2
\]
that expresses the acute angle the trail makes with the horizontal at the point $s$ miles from the trailhead in terms of the acute angle the curve $h=f(s)$ makes with the $s$-axis at the point $(s,f(s))$. 

\item Evaluate the derivative
\[
\frac{d\phi}{d\theta}\Big|_{\theta = \arctan(2)} .
\]
\item Interpret the derivative in part (b) as a scaling factor. Give an example with specific small changes.  
\end{enumerate}
\end{question}



\begin{question} \label{Qpdfe9930}
The function
\[
      T = f(r) = \frac{R^{3/2}}{\sqrt{2GM}} \left( \frac{\pi}{2} - \arcsin\left( \sqrt{\frac{r}{R}} \right) + \sqrt{\frac{r}{R}\left(1 - \frac{r}{R} \right)}    \right)
\]

\end{question}


\end{document}